\documentclass[exercice,a4]{thedocument}%
\usepackage{pstricks}%
\usepackage{enumitem}%
\usepackage{graphicx}%
\startdatakeys
    \source{25GENMATAS1}
    \BO{2026}
    \cycle{4}
    \competencesDuSocle{CA,RA}
    \objectifsApprentissage{C4A1aa,C4A2ad,C4B2cb,C4D2aa,C4D3ac}
    \motsCles{volume,pavé droit,cylindre,aire,conversion}
\enddatakeys
\begin{document}%
\providecommand{\degres}{^\circ}%
    

\emph{Dans cet exercice, les deux parties sont indépendantes}

\subsection*{Rappels}
\begin{itemize}[label=$\bullet~~$]
\item Volume du cylindre = Aire de la base $\times$ Hauteur du cylindre
\item Aire du disque = $\pi \times$ (rayon)$^2$
\item $1$\,cm$^3 = 1$\,mL
\end{itemize}

Pour un anniversaire, on veut préparer des cocktails de jus de fruits.

\subsection*{Partie 1 : Étude des glaçons}

\textbf{Document} : photo du moule à glaçons utilisé et caractéristiques des glaçons :

\medskip

\begin{minipage}{0.48\linewidth}
\begin{center}

\eps[width=6cm]{moule1}
\end{center}
\end{minipage}\hfill
\begin{minipage}{0.48\linewidth}
Chaque glaçon a la forme d'un pavé droit :
\begin{itemize}[label=$\bullet~~$]
\item de longueur 5 cm ;
\item de largeur 2,5 cm ;
\item de hauteur 1,5 cm.
\end{itemize}
\end{minipage}
\begin{enumerate}
\item On possède 12 moules à glaçons de ce type. Combien peut-on faire de glaçons en même temps ?
\item Montrer que le volume d'un glaçon est d'environ $19$~mL.
\item 5 litres d'eau sont-ils suffisants pour remplir ces 12 moules à glaçons ?
\end{enumerate}

\subsection*{Partie 2 : Le service}

\begin{minipage}{0.68\linewidth}
On souhaite servir le cocktail dans des verres cylindriques.

\begin{enumerate}[resume]
\item Montrer que le verre a un volume total d'environ $295$~mL.
\item Pour verser précisément $25$ cL de cocktail, on utilise des verres avec un repère indiquant une contenance de $25$~cL.
	\begin{enumerate}
		\item On a préparé $30$ litres de cocktail. Combien peut-on remplir de verres contenant $25$ cL de cocktail ?
		\item En versant $25$ cL de cocktail dans le verre, à quelle hauteur $h$ du verre, le liquide arrive-t-il ?

Arrondir au dixième.
	\end{enumerate}
\end{enumerate}
\end{minipage}\hfill
\begin{minipage}{0.3\linewidth}
\psset{unit=1cm,arrowsize=2pt 3}
\begin{pspicture}(-2.7,0)(2.7,7.7)
\psellipse[linewidth=1.5pt](0,1.4)(1.05,0.3)
\psellipse[linewidth=1.5pt](0,1.45)(1.05,0.3)
\psellipse[linewidth=1.5pt](0,6.6)(1.05,0.3)
\psellipse[linewidth=1.5pt](0,6.65)(1.05,0.3)
\psline[linewidth=1.5pt](-1.05,1.4)(-1.05,6.65)
\psline[linewidth=1.5pt](1.05,1.4)(1.05,6.65)
\psellipse[linewidth=1.5pt,linestyle=dotted](0,5.7)(1.05,0.3)
\rput(0,7.6){Caractéristiques d'un verre}
\rput(1.8,6.3){Repère}\psline[linewidth=0.9pt]{->}(1.3,6.2)(0.6,5.44)
\psline[linewidth=0.9pt]{<->}(-1.5,1.4)(-1.5,6.6)\uput[l](-1.5,4){15 cm}
\psline[linewidth=0.9pt]{<->}(1.3,1.4)(1.3,5.6)\uput[r](1.3,3.5){$h$}
\psline[linewidth=0.9pt]{<->}(-1.05,0.8)(1.05,0.8)\uput[d](0,0.8){5 cm}
\end{pspicture}
\end{minipage}
\showthesource
\end{document}
